% GENERATED FILE. Edit canonical lessons, then run npm run generate:books.
% canonical-source-hash: fnv1a64:b3dd938b100a414c
% canonical-lessons: MR-C264-vrddha, MR-C264-jule, MR-C264-sasu, MR-C264-sasre, MR-C264-javai

\chapter{Elderly person, Twins, Mother-in-law, Father-in-law, Son-in-law}
\label{ch:mr-elderly-person-twins-mother-in-law-father-in-law-son-in-law}
\hlchaptermodality{264}

\begin{chapteropening}
\textbf{By the end of this chapter:} \emph{I can say an elderly person, twins, a mother-in-law, a father-in-law and a son-in-law.}

Complete the last lesson of chapter 264: I can say an elderly person, twins, a mother-in-law, a father-in-law and a son-in-law.
\end{chapteropening}

\section[\texorpdfstring{\mr{वृद्ध} (vṛddha)}{वृद्ध (vṛddha)}]{\mr{वृद्ध} (vṛddha) — an elderly person}

\label{lesson:MR-C264-vrddha}



\begin{quote}
Before the new one: say the Marathi for an actress, then the Marathi for a poem.
\end{quote}

\subsection*{You'll want to know: \mr{वृद्ध}}
\textbf{\mr{वृद्ध}} — \emph{vṛddha} — \textquotedblleft{}an elderly person\textquotedblright{}.

\begin{grammarlens}[title={What you've built}]
One more word for this chapter.
\end{grammarlens}

\subsection*{Your turn}
\begin{itemize}
\raggedright
  \item \emph{Say it:} \emph{vṛddha}
  \item \emph{Say it:} \emph{vṛddha}, once more
  \item \emph{Say it:} say \emph{vṛddha}
  \item \emph{Recall:} say the Marathi for an actress, then the Marathi for a poem, then say \emph{vṛddha} again
\end{itemize}

\begin{tcolorbox}[breakable,colback=teal!4,colframe=teal!35!black,title={Before you move on}]
What does \mr{वृद्ध} mean? (\textquotedblleft{}An elderly person\textquotedblright{}.) Say it once more.
\end{tcolorbox}
\section[\texorpdfstring{\mr{जुळे} (juḷe)}{जुळे (juḷe)}]{\mr{जुळे} (juḷe) — twins}

\label{lesson:MR-C264-jule}



\begin{quote}
Before the new one: say the Marathi for a poem, then the Marathi for an elderly person.
\end{quote}

\subsection*{You'll want to know: \mr{जुळे}}
\textbf{\mr{जुळे}} — \emph{juḷe} — \textquotedblleft{}twins\textquotedblright{}.

\begin{grammarlens}[title={What you've built}]
One more word for this chapter.
\end{grammarlens}

\subsection*{Your turn}
\begin{itemize}
\raggedright
  \item \emph{Say it:} \emph{juḷe}
  \item \emph{Say it:} \emph{juḷe}, once more
  \item \emph{Say it:} say \emph{juḷe}
  \item \emph{Recall:} say the Marathi for a poem, then the Marathi for an elderly person, then say \emph{juḷe} again
\end{itemize}

\begin{tcolorbox}[breakable,colback=teal!4,colframe=teal!35!black,title={Before you move on}]
What does \mr{जुळे} mean? (\textquotedblleft{}Twins\textquotedblright{}.) Say it once more.
\end{tcolorbox}
\section[\texorpdfstring{\mr{सासू} (sāsū)}{सासू (sāsū)}]{\mr{सासू} (sāsū) — a mother-in-law}

\label{lesson:MR-C264-sasu}



\begin{quote}
Before the new one: say the Marathi for an elderly person, then the Marathi for twins.
\end{quote}

\subsection*{You'll want to know: \mr{सासू}}
\textbf{\mr{सासू}} — \emph{sāsū} — \textquotedblleft{}a mother-in-law\textquotedblright{}.

\begin{grammarlens}[title={What you've built}]
One more word for this chapter.
\end{grammarlens}

\subsection*{Your turn}
\begin{itemize}
\raggedright
  \item \emph{Say it:} \emph{sāsū}
  \item \emph{Say it:} \emph{sāsū}, once more
  \item \emph{Say it:} say \emph{sāsū}
  \item \emph{Recall:} say the Marathi for an elderly person, then the Marathi for twins, then say \emph{sāsū} again
\end{itemize}

\begin{tcolorbox}[breakable,colback=teal!4,colframe=teal!35!black,title={Before you move on}]
What does \mr{सासू} mean? (\textquotedblleft{}A mother-in-law\textquotedblright{}.) Say it once more.
\end{tcolorbox}
\section[\texorpdfstring{\mr{सासरे} (sāsre)}{सासरे (sāsre)}]{\mr{सासरे} (sāsre) — a father-in-law}

\label{lesson:MR-C264-sasre}



\begin{quote}
Before the new one: say the Marathi for twins, then the Marathi for a mother-in-law.
\end{quote}

\subsection*{You'll want to know: \mr{सासरे}}
\textbf{\mr{सासरे}} — \emph{sāsre} — \textquotedblleft{}a father-in-law\textquotedblright{}.

\begin{grammarlens}[title={What you've built}]
One more word for this chapter.
\end{grammarlens}

\subsection*{Your turn}
\begin{itemize}
\raggedright
  \item \emph{Say it:} \emph{sāsre}
  \item \emph{Say it:} \emph{sāsre}, once more
  \item \emph{Say it:} say \emph{sāsre}
  \item \emph{Recall:} say the Marathi for twins, then the Marathi for a mother-in-law, then say \emph{sāsre} again
\end{itemize}

\begin{tcolorbox}[breakable,colback=teal!4,colframe=teal!35!black,title={Before you move on}]
What does \mr{सासरे} mean? (\textquotedblleft{}A father-in-law\textquotedblright{}.) Say it once more.
\end{tcolorbox}
\section[\texorpdfstring{\mr{जावई} (jāvaī)}{जावई (jāvaī)}]{\mr{जावई} (jāvaī) — a son-in-law}

\label{lesson:MR-C264-javai}



\begin{quote}
Before the new one: say the Marathi for a mother-in-law, then the Marathi for a father-in-law.
\end{quote}

\subsection*{You'll want to know: \mr{जावई}}
\textbf{\mr{जावई}} — \emph{jāvaī} — \textquotedblleft{}a son-in-law\textquotedblright{}.

\begin{grammarlens}[title={What you've built}]
One more word for this chapter.
\end{grammarlens}

\subsection*{Your turn}
\begin{itemize}
\raggedright
  \item \emph{Say it:} \emph{jāvaī}
  \item \emph{Say it:} \emph{jāvaī}, once more
  \item \emph{Say it:} say \emph{jāvaī}
  \item \emph{Recall:} say the Marathi for a mother-in-law, then the Marathi for a father-in-law, then say \emph{jāvaī} again
\end{itemize}

\begin{tcolorbox}[breakable,colback=teal!4,colframe=teal!35!black,title={Before you move on}]
What does \mr{जावई} mean? (\textquotedblleft{}A son-in-law\textquotedblright{}.) Say it once more.
\end{tcolorbox}
